\section{Consequences for controlled validation}
\label{sec:discussion}
An initialization audit should precede attributing every later difference to
the heat-flux closure. The current test exposes a mismatch in physical moments
despite matching algebraic drift expressions. A controlled repeat should match
density, ordinary mean velocity and the selected pressure definition, and
report the initial magnetic curl and moment current on their actual staggered
locations. If relativistic particle initialization and a nonrelativistic fluid
model are compared, the conversion and the target moment convention must be
stated explicitly. The cold proper-velocity conversion alone is not a proof
that finite-temperature moments have been matched.

The clean one-fluid endpoint contrast establishes a numerical response to
the aggregate-moment intervention: all four tested endpoint errors decrease.
It does not establish a physical cause of the source's reconnection differences.
The intervention replaces the conserved moment vector through second order,
including sampled spatial variation; it is not confined to the mean current.
Both one-fluid saved histories have covariance
failures, the clean unmatched two-fluid endpoint is rejected, and the matched
two-fluid run fails before producing its comparison endpoint. A smaller
endpoint error cannot erase those failures or substitute for realizability
and resolution checks. We therefore do not infer that correcting the drift
eliminates the reported density deficit, reproduces full PIC heat flux with
one fluid, or reverses the source's physical conclusion.

The archived one- and two-fluid endpoints remain positive definite under
the same tensor test. This distinguishes an observed failure of the identified
clean implementation from a claim about the unavailable dirty executable.
The finite snapshot archive does not establish that the original intermediate
trajectory was admissible, either. Identified executable provenance, full-tensor
state checks and grid-converged controls are needed before transferring the
numerical endpoint changes into a physical causal statement.

The auxiliary homogeneous stability test supplies a separate calibration.
Retaining population drift creates the possibility of streaming modes; it
does not ensure the correct kinetic susceptibility or growth rate. The
implemented relaxation should be included in the dispersion calculation.
Finite relaxation can change a spectrum even when the unperturbed pressure
is isotropic and the relaxation source vanishes at equilibrium. In the stated
example the resulting growing mode is a false positive relative to the
collisionless kinetic reference. This is a test of a model's continuum
response, not evidence that the original simulation visibly contains that mode.
The new clean-source evolutions do not provide a source-sheet eigenanalysis
or a physically validated, grid-converged matched comparison. This diagnostic
comment quantifies numerical sensitivity and exposes validation requirements;
the source's nonlinear physical conclusions are not thereby disproved.

The pressure-source diagnostic should also distinguish a third-moment amplitude
from its divergence. The scale proxy $p\,d_e/|q|$ used in the source is a useful
estimate when $d_e$ is an appropriate gradient length. It is not identical to
$p/|\nabla\!\cdot q|$. A spatially constant nonzero heat-flux tensor has a finite
amplitude proxy and zero divergence. The source explicitly notes that its
isotropization closure constrains a divergence and omits the reconstructed
fluid tensor from some plots; we do not characterize that disclosure as a
hidden omission. Direct source-term comparisons and rate convergence would
make the approximation more informative.

Finally, the source's time-maximum spatial score is a best-time diagnostic,
not a uniform-in-time error bound or a forward prediction test at every scan
point. The conditioned estimate above explains why both absolute residuals
and a nondegenerate denominator are needed when assessing relative accuracy.
These qualifications preserve the observed scan and suggest measurements
needed before extending its interpretation to new current sheets, later
island recirculation, or turbulence.

The valid component-to-aggregate moment decomposition remains the foundation
of this analysis. Additional labels carry information absent from a generic
state described only by total density, velocity and pressure. The corrections
identify controls and assumptions needed for using that additional information
reliably, rather than a general rejection of multicomponent fluid modeling.
